LW provides a uniquely intuitive compactness metric for courtroom communication.

\subsection{When to Use LW}

\begin{itemize}
  \item \textbf{Primary use}: When the legal challenge involves ``elongation''
    claims---districts that run as corridors across the state.
  \item \textbf{Secondary use}: As a complement to Reock in the composite profile,
    because both measure elongation but from different geometric primitives
    (MBC vs.\ MBR).
  \item \textbf{Threshold certification}: When the jurisdiction applies a
    LW $> 3$ threshold (Illinois, New York expert practice), use LW as the
    primary pass-fail metric.
\end{itemize}

\subsection{Courtroom Explanation}

\begin{quote}
  \emph{The Length-Width Ratio measures the aspect ratio of the smallest rectangle
  that can contain the district, oriented at any angle. A district with LW $= 3$
  is three times as long as it is wide. All 14 NC congressional districts in
  the proposed plan have LW $\leq 2.38$, well below the 3:1 threshold that
  courts have applied in redistricting challenges.}
\end{quote}

\subsection{Disclosure Statement}

For court filings:
\begin{quote}
  \emph{All Length-Width Ratio scores were computed using the minimum-area
  bounding rectangle (MBR) of district convex hulls, projected to Albers
  Equal Area Conic (EPSG:5070). The MBR was computed using the rotating
  calipers algorithm applied to the convex hull vertices. LW $= \ell / w$
  where $\ell \geq w$ are the MBR dimensions. Axis-aligned bounding boxes
  (AABB) were not used because they are orientation-dependent.}
\end{quote}

\subsection{CLI Usage}

\begin{verbatim}
bisect label-analyze <label> --year 2020 --types lw
\end{verbatim}
Output includes both MBR and AABB values for comparison, with a flag
if any district exceeds LW $= 3.0$.
